\section{Finite-data testing on the original behavior image}
\label{sec:v18-audit}
\subsection{The reset interface}
Let $B\subset\prod_{j=1}^J\Delta(\Omega_j)$ be compact, and let
$Q=(Q_j)$ be fixed. Set $d_j=|\Omega_j|$, $D=\sum_j(d_j-1)$,
$\mathcal I=\{(j,A):A\subseteq\Omega_j\}$, and
$f_{j,A}(b)=Q_j(A)-b_j(A)$. Include the empty event. No convexity of $B$
is assumed.

A candidate $b$ is fixed for the whole audit. The auditor obtains $n$
independent samples of each $b_j$, with independence between rows.
In a causal experiment, row $j$ means a specified parameter and deterministic
feedback tree; the reset interface supplies that row. It is an additional
experimental interface, not information available to the online simulator.
A complete sample contains the terminal mark. Private randomness and any
per-run hidden selector are reset between runs. The design of the candidate
is not resampled.

Write $C_{j,w}$ for the counts, and $\mathcal N_n$ for all count arrays
with $\sum_wC_{j,w}=n$. Their probability is the polynomial
\[
 p_C(b)=\prod_j\frac{n!}{\prod_w C_{j,w}!}
                    \prod_w b_j(w)^{C_{j,w}}.
\]
A randomized audit is an array $c=(c_{C,i})$ with $c_C\in\Delta(\mathcal I)$.
After training it chooses $i=(j,A)$ according to $c_C$. Independent candidate
and target validation runs give the signed score
$\1_A(\hbox{target})-\1_A(\hbox{candidate})$. Its expected value is
\begin{equation}\label{eq:v18-auditpayoff}
 g_c(b)=\sum_{C,i}p_C(b)c_{C,i}f_i(b),\qquad
 \vaudit_n(B,Q)=\max_c\min_{b\in B}g_c(b).
\end{equation}
Thus every audit payoff is polynomial of total degree at most $Jn+1$ in
the behavior coordinates. The probability of a large observed score in a
single audit is not asserted to equal its expectation; additional
independent validation would be required for a specified confidence level.

\begin{theorem}[Reset-audit duality and finite-sample approximation]
\label{thm:v18-audit}
Let $\delta=\min_{b\in B}\max_i f_i(b)$. Then all the extrema below are
attained, and
\begin{align}
 \vaudit_n(B,Q)
 &=\min_{\mu\in\mathcal P(B)}
       \sum_{C\in\mathcal N_n}\max_{i\in\mathcal I}
                  \int_B p_C(b)f_i(b)\,d\mu(b),
                    \label{eq:v18-auditdual}\\
 0&\le\delta-\vaudit_n(B,Q)\le\sqrt{D/n}.
                    \label{eq:v18-auditrate}
\end{align}
The values are nondecreasing in $n$. If $a=\dim\aff B$, a minimizing
measure in \eqref{eq:v18-auditdual} can be supported on at most
$\binom{a+Jn+1}{a}$ points.

An explicit feasible audit attaining the lower bound
$\delta-\sqrt{D/n}$ chooses an event maximizing
$Q_j(A)-C_j(A)/n$, with a fixed deterministic tie rule. If the rational
target probabilities are given, this is a rationally specified finite
rule. With a free schedule clock, a conservative retained-state bound is
\[
 H(|\mathcal I|+1)(n+1)^{\sum_jd_j},
\]
where $H$ is the number of states needed to execute and record one causal
transcript. This budget is charged to the auditor, not to the simulator.
\end{theorem}
\begin{proof}
The coefficient arrays form a compact convex product of finite simplexes.
The integrand is continuous in $b$ and affine in $c$. The probability
measures on compact $B$ form a compact convex set for weak convergence.
The minimum over $b$ equals the minimum over such measures for fixed $c$.
Finite-dimensional minimax (or its compact measure version) therefore
exchanges the two extrema. Maximization over each independent coefficient
simplex selects its largest coordinate integral, giving
\eqref{eq:v18-auditdual}. Continuity and compactness give attainment.
This step uses measures on the original $B$, not feasible private
realizations of their barycenters.

For any audit, $g_c(b)\le F(b)=\max_i f_i(b)$. At a minimizer of $F$ this
gives $\vaudit_n\le\delta$. The zero event gives $\vaudit_n\ge0$.
For the stated explicit rule put $\widehat b_j=C_j/n$ and
$\rho=\max_j\|\widehat b_j-b_j\|_{\rm TV}$. Comparing the empirical
maximizer with a true maximizing event gives the pointwise regret bound
\[
 F(b)-f_{i(C)}(b)\le2\rho.
\]
For a multinomial empirical measure,
\[
 \E\|\widehat b_j-b_j\|_{\rm TV}^2
 \le\frac{d_j}{4}\sum_w\E(\widehat b_j(w)-b_j(w))^2
 =\frac{d_j(1-\sum_wb_j(w)^2)}{4n}
 \le\frac{d_j-1}{4n}.
\]
Consequently $\E\rho\le(\sum_j\E\|\widehat b_j-b_j\|_{\rm TV}^2)^{1/2}
\le\sqrt{D/n}/2$. Averaging the regret uniformly over $B$ proves
\eqref{eq:v18-auditrate}.

From $n+1$ unordered samples in each row, discard one uniformly at random.
The remaining counts have exactly the $n$-sample law, independently across
rows. This realizes every $n$-sample audit as a randomized $(n+1)$-sample
audit and proves monotonicity.

Choose affine coordinates on $\aff B$. All the functions $p_C f_i$ lie
in the span of monomials of total degree at most $Jn+1$, a space of
dimension at most $s=\binom{a+Jn+1}{a}$ including the constant. The moment
vector of a probability measure is in the convex hull of the evaluation
vectors; their constant coordinate is one, so their affine dimension is
at most $s-1$. Eliminating affine dependencies gives a representation by
at most $s$ points. It preserves every integral in
\eqref{eq:v18-auditdual}, and hence preserves the objective.

Finally, each count is in $\{0,\ldots,n\}$. The indicated product register
stores all counts, the within-run state, and the selected event for
validation. Only the latter event and the within-run buffer need survive
once training ends, so the product bound is conservative. Randomizing an
event uses fresh randomness. The schedule, reset access and terminal mark
readout were declared before the theorem. No private simulator state is
read by this implementation.
\end{proof}

The dual measure chooses a candidate once, after which its samples are
independent conditional on that candidate. Replacing it by
$\bar b=\int b\,d\mu$ would replace $\int p_C(b)f_i(b)d\mu$ by
$p_C(\bar b)f_i(\bar b)$, generally a different quantity. It would also
replace the joint replicated law by a product of its marginals. Neither
operation is part of the theorem. This is the operational meaning of the
moment side, rather than a same-budget convexification theorem.

\begin{proposition}[Both sides sampled]
\label{prop:v18-paired}
Suppose the target reference $Q^b$ varies with the candidate, as it does
when visible-selector weights are design variables. Supply independent
training samples from both $b$ and $Q^b$, and use their empirical
difference. If their row dimensions are $D_P,D_Q$, the same construction
has regret at most
$\sqrt{D_P/n}+\sqrt{D_Q/n}$. Minimax and attainment hold on the compact
image of pairs $(b,Q^b)$. In particular a variable reference law must not
be treated as known fixed input to Theorem~\ref{thm:v18-audit}.
\end{proposition}
\begin{proof}
The error in any event discrepancy is at most the sum of the maximal
row variations of the two empirical tables. Apply the preceding second
moment bound separately and then the empirical-maximizer regret inequality.
Count probabilities are products for the two tables. They are continuous
polynomials on the compact image of pairs, so the same minimax argument
applies without changing the simulator signature.
\end{proof}

\subsection{A matching uniform sampling obstruction}
\begin{theorem}[Binary reset lower bound]
\label{thm:v18-auditlower}
Take one binary row with target $\operatorname{Ber}(1/2)$ and candidate
class $B_h=\{\operatorname{Ber}(1/2-h),\operatorname{Ber}(1/2+h)\}$,
where $0<h<1/2$. Then
\begin{equation}\label{eq:v18-binaryaudit}
 \delta=h,\qquad
 \vaudit_n=h\left\|\operatorname{Ber}(1/2+h)^{\otimes n}
               -\operatorname{Ber}(1/2-h)^{\otimes n}\right\|_{\rm TV}.
\end{equation}
For every $n\ge1$, choosing $h=1/(8\sqrt n)$ gives
\[
 \delta-\vaudit_n\ge\frac1{16\sqrt n}.
\]
Thus the exponent in \eqref{eq:v18-auditrate} is sharp uniformly over
compact candidate classes, already at affine dimension one. This does
not assert a degree lower bound for revealed-behavior polynomials, nor
that this two-point class equals an unrestricted private-row image.
\end{theorem}
\begin{proof}
Only the two complementary singleton events and the zero event matter.
A randomized choice is a statistic $d$ in $[-1,1]$. The two expected
payoffs are $h\E_+d$ and $-h\E_-d$. Their minimum is at most half their
sum, which is at most $h\|P_+^n-P_-^n\|_{\rm TV}$. Conversely take the
sign of the likelihood difference, assigning zero at a tie. Complementing
every sample interchanges $P_+^n$ and $P_-^n$ and changes that sign.
The two payoffs are equal, and their average attains the variation bound.
The statistic depends only on the count, so it is an admissible audit.
The symmetric measure on the two candidates attains the dual value.
Every Dirac measure has dual objective $h$, whereas for finite $n$ the two
product laws are equivalent and their variation is strictly less than one.
Thus no dual minimizer is a Dirac measure. Their barycenter is the target
law, but replacing the dual by that barycenter would incorrectly give
value zero instead of the strictly positive value in
\eqref{eq:v18-binaryaudit}.

For $h\le1/8$, elementary integration of $(1-u^2)^{-1}$ gives
\[
 \KL(\operatorname{Ber}(1/2+h)\,\|\,\operatorname{Ber}(1/2-h))
 =2h\log\frac{1+2h}{1-2h}\le16h^2.
\]
The product relative entropy is at most $16nh^2$. Pinsker's inequality
bounds its total variation by $\sqrt{8nh^2}$. At the specified $h$ this
is less than $1/2$. Substituting in \eqref{eq:v18-binaryaudit} proves the
claim. The inequality uses the classical two-point testing mechanism;
the restriction being measured is the reset budget of the present game.
\end{proof}

\subsection{Transport of one fixed audit}
\begin{proposition}[Reset cost under model perturbation]
\label{prop:v18-audittransport}
Let $q$ range over the same compact set of executable row designs in two
models. Suppose uniformly in $q,j$ that
\[
 \|P_j^q-\widetilde P_j^q\|_{\rm TV}\le\alpha,\qquad
 \|Q_j-\widetilde Q_j\|_{\rm TV}\le\beta.
\]
For the same coefficient rule $c$, evaluated with each model's own target
validation,
\begin{equation}\label{eq:v18-audittransport}
 |g_c(q)-\widetilde g_c(q)|\le(2Jn+1)\alpha+\beta.
\end{equation}
The deficiencies themselves differ by at most $\alpha+\beta$, at the
same resource budget. If the source behavior is unchanged, a fixed audit
loses at most $\beta$, independently of its reset count.
\end{proposition}
\begin{proof}
A product coupling of the $Jn$ training runs fails with probability at
most $Jn\alpha$. Taking counts and choosing the event are contractions.
A score with fixed target discrepancies lies in $[-1,1]$, so the change
in its expectation from this training-law perturbation is at most
$2Jn\alpha$. At each event the discrepancy changes by at most
$\alpha+\beta$. This proves \eqref{eq:v18-audittransport}. For deficiency,
every row variation changes by at most $\alpha+\beta$; taking the maximum
and then the minimum preserves that bound. No new simulator is inserted.
\end{proof}

For a rational polynomial image, a rational rule and a rational threshold,
strict positivity of $g_c(q)$ above that threshold is a first-order formula
over the original row simplexes. Real quantifier elimination decides it.
The same is true of an existential rule followed by this universal
condition. This is an exact finite verification statement, not an efficient
algorithm. The explicit empirical rule below will be verified by a uniform
analytic inequality, rather than by storing its enormous complete table.
